\begin{thebibliography}{00}
\bibitem{Lewis2020} P. Lewis et al., ``Retrieval-Augmented Generation for Knowledge-Intensive NLP Tasks,'' in \textit{Advances in Neural Information Processing Systems}, 2020.
\bibitem{Gao2023} Y. Gao et al., ``Retrieval-Augmented Generation for Large Language Models: A Survey,'' arXiv:2312.10997, 2023.
\bibitem{Edge2024} D. Edge et al., ``From Local to Global: A Graph RAG Approach to Query-Focused Summarization,'' arXiv:2404.16130, 2024.
\bibitem{Guo2024} Z. Guo et al., ``LightRAG: Simple and Fast Retrieval-Augmented Generation,'' arXiv:2410.05779, 2024.
\bibitem{Gutierrez2024} J. Guti\'errez et al., ``HippoRAG: Neurobiologically Inspired Long-Term Memory for Large Language Models,'' arXiv:2405.14831, 2024.
\bibitem{Li2025TGRAG} D. Li et al., ``T-GRAG: A Dynamic GraphRAG Framework for Resolving Temporal Conflicts and Redundancy in Knowledge Retrieval,'' arXiv:2508.01680, 2025.
\bibitem{Han2025TGRAG} J. Han et al., ``RAG Meets Temporal Graphs: Time-Sensitive Modeling and Retrieval for Evolving Knowledge,'' arXiv:2510.13590, 2025.
\bibitem{Xiang2025} Z. Xiang et al., ``When to Use Graphs in RAG: A Comprehensive Analysis for Graph Retrieval-Augmented Generation,'' arXiv:2506.05690, 2025.
\end{thebibliography}
